\documentclass{book}
\usepackage[backend=biber,natbib=true,style=authoryear]{biblatex}
\addbibresource{/home/nguyen/1_NQBH/math/bib.bib}
\usepackage[utf8]{inputenc}
\usepackage{float}
\usepackage{graphicx}
\usepackage[colorlinks=true,linkcolor=blue,urlcolor=red,citecolor=magenta]{hyperref}
\usepackage{amsmath,amssymb,amsthm}
\allowdisplaybreaks
\numberwithin{equation}{section}
\newtheorem{assumption}{Assumption}[section]
\newtheorem{lemma}{Lemma}[section]
\newtheorem{corollary}{Corollary}[section]
\newtheorem{definition}{Definition}[section]
\newtheorem{proposition}{Proposition}[section]
\newtheorem{theorem}{Theorem}[section]
\newtheorem{notation}{Notation}[section]
\newtheorem{remark}{Remark}[section]
\newtheorem{example}{Example}[section]
\newtheorem{ques}{Question}[section]
\newtheorem{problem}{Problem}[section]
\newtheorem{conjecture}{Conjecture}[section]
\usepackage[left=0.5in,right=0.5in,top=1.5cm,bottom=1.5cm]{geometry}
\usepackage{fancyhdr}
\pagestyle{fancy}
\fancyhf{}
\addtolength{\headheight}{0pt} % obsolete
\lhead{\scriptsize \chaptername~\thechapter}
\rhead{\itshape \scriptsize \nouppercase{\leftmark}} %\nouppercase !
\renewcommand{\chaptermark}[1]{\markboth{#1}{}}
\cfoot{\thepage}

\title{G. H. Hardy (2017). \textit{A Mathematician's Apology}}
\author{Collector: Nguyen Quan Ba Hong}
\date{\today}

\begin{document}
\maketitle
\setcounter{secnumdepth}{6}
\setcounter{tocdepth}{6}
\tableofcontents

%------------------------------------------------------------------------------%

\chapter*{Foreword}

``It was a perfectly ordinary night at Christ's high table, except that Hardy was dining as a guest.

He had just returned to Cambridge as Sadleirian professor, and I had heard something of him from young Cambridge mathematicians.

They were delighted to have him back:\fbox{ he was a \textit{real} mathematician}, they said, not like those Diracs and Bohrs the physicists were always talking about: \fbox{\textit{he was the purest of the pure}.}

He was also unorthodox, eccentric, radical, ready to talk about anything.

This was 1931, and the phrase was not yet in English use, but in later days they would have said that in some indefinable way he had star quality.

%
So, from lower down the table, I kept studying him.

He was then in his early 50s: his hair was already grey, above skin so deeply sunburnt that it stayed a kind of Red Indian bronze.

His face was beautiful - high cheek bones, thin nose, spiritual and austere but capable of dissolving into convulsions of internal gamin-like amusement.

He had opaque brown eyes, bright as a bird's - a kind of eye not uncommon among those with a gift for conceptual thought.

Cambridge at that time was full of unusual and distinguished faces - but even then, I thought that night, Hardy's stood out.

%
I do not remember what he was wearing.

It may easily have been a sports coat and grey flannels under his gown.

Like Einstein, he dressed to please himself: though, unlike Einstein, he diversified his casual clothing by a taste for expensive silk shirts.

%
As we sat round the combination-room table, drinking wine after dinner, someone said that Hardy wanted to talk to me about cricket.

I had been elected only a year before, but Christ's was then a small college, and the pastimes of even the junior fellows were soon identified.

I was taken to sit by him.

I was not introduced.

He was, as I later discovered, shy and self-conscious in all formal actions, and had a dread of introductions.

He just put his head down as it were in a butt of acknowledgment, and without any preamble whatever began:
\begin{quotation}
    `You're supposed to know something about cricket, aren't you?' Yes, I said, I knew a bit.
\end{quotation}
Immediately he began to put me through a moderately stiff viva.

Did I play?

What sort of performer was I?

I half-guessed that he had a horror of persons, then prevalent in academic society, who devotedly studied the literature but had never played the game.

I trotted out my credentials, such as they were.

He appeared to find the reply partially reassuring, and went on to more tactical questions.

Whom should I have chosen as captain for the last test match a year before (in 1930)?

If the selectors had decided that Snow was the man to save England, what would have been my strategy and tactics?

(`You are allowed to act, if you are sufficiently modest, as non-playing captain.')

And so on, oblivious to the rest of the table.

He was quite absorbed.

%
As I had plenty of opportunities to realize in the future, Hardy had no faith in intuitions or impressions, his own or anyone else's.

The only way to assess someone's knowledge, in Hardy's view, was to examine him.

That went for mathematics, literature, philosophy, politics, anything you like.

If the man had bluffed and then wilted under the questions, that was his lookout.

1st things came 1st, in that brilliant and concentrated mind.

%
That night in the combination-room, it was necessary to discover whether I should be tolerable as a cricket companion.

Nothing else mattered.

In the end he smiled with immense charm, with child-like openness, and said that Fenner's (the university cricket ground) next season might be bearable after all, with the prospect of some reasonable conversation.

%
Thus, just Is I owed my acquaintanceship with Lloyd George to his passion for phrenology, I owed my friendship with Hardy to having wasted a disproportionate amount of my youth on cricket.

I don't know what the moral is.

But it was a major piece of luck for me.

This was intellectually the most valuable friendship of my life.

His mind, as I have just mentioned, was brilliant and concentrated: so much so that by his side anyone else's seemed a little muddy, a little pedestrian and confused.

He wasn't a great genius, as Einstein and Rutherford were.

He said, with his usual clarity, that if the word meant anything he was not a genius at all.

At his best, he said, he was for a short time the 5th best pure mathematician in the world.

Since his character was as beautiful and candid as his mind, he always made the point that his friend and collaborator Littlewood was an appreciably more powerful mathematician than he was, and that his prot\'eg\'e Ramanujan really had natural genius in the sense (though not to the extent, and nothing like so effectively) that the greatest mathematicians had it.

%
People sometimes thought he was under-rating himself, when he spoke of these friends.

It is true that he was magnanimous, as far from envy as a man can be: but I think one mistakes his quality if one doesn't accept his judgment.

I prefer to believe in his own statement in \textit{A Mathematician Apology}, at the same time so proud and so humble:
\begin{quotation}
    `I still say to myself when I am depressed and find myself forced to listen to pompous and tiresome people, ``Well, I have done 1 thing you could never have done, and that is to have collaborated with Littlewood and Ramanujan on something like equal terms.''
\end{quotation}
In any case, his precise ranking must be left to the historians of mathematics (though it will be an almost impossible job, since so much of his best work was done in collaboration).

There is something else, though, at which he was clearly superior to Einstein or Rutherford or any other great genius: and that is at \textit{turning any work of the intellect, major or minor or sheer play, into a work of art}.

It was that gift above all, I think, which made him, almost without realizing it, purvey such intellectual delight.

When \textit{A Mathematician's Apology} was 1st published, Graham Greene in a review wrote that long with Henry James's notebooks, this was the best account of what it was like to be a \textit{creative artist}.

Thinking about the effect Hardy had on all those around him, I believe that is the clue.

%
He was born, in 1877, into a modest professional family.

His father was Bursar and Art Master at Cranleigh, then a minor public (English for private) school.

His mother had been senior mistress at the Lincoln Training College for teachers.

Both were gifted and mathematically inclined.

In his case, as in that of most mathematicians, the gene pool doesn't need searching for.

Much of his childhood, unlike Einstein's, was typical of a future mathematician's.

He was demonstrating a formidably high I.Q. as soon as, or before, he learned to talk.

At the age of 2 he was writing down numbers up to millions (a common sign of mathematical ability).

When he was taken to church he amused himself by factorizing the numbers of the hymns: he played with numbers from that time on, a habit which led to the touching scene at Ramanujan's sick-bed: the scene is well known, but later on I shall not be able to resist repeating it.

%
It was an enlightened, cultivated, highly literate Victorian childhood.

His parents were probably a little obsessive, but also very kind.

Childhood in such a Victorian family was as gentle a time as anything we could provide, though probably intellectually somewhat more exacting.

His was unusual in just 2 respects.

In the 1st place, he suffered from an acute self-consciousness at an unusually early age, long before he was 12.

His parents knew he was prodigiously clear, and so did he.

He came top of his class in all subjects.

But, as the result of coming top of his class, he had to go in front of the school to receive prizes: and that he could not bear.

Dining with me 1 night, he said that he deliberately used to try to get his answers wrong so as to be spared this intolerable ordeal.

His capacity for dissimulation, though, was always minimal: he got the prizes all the same.

%
Some of this self-consciousness wore off.

He became competitive.

As he says in the \textit{Apology}:
\begin{quotation}
    `I do not remember having felt, as a boy, any \textit{passion} for mathematics, and such notions as I may have had of the career of a mathematician were far from noble.
    
    I thought of mathematics in terms of examinations and scholarships: I wanted to beat other boys, and this seemed to be the way in which I could do so most decisively.'
\end{quotation}
Nevertheless, he had to live with an over-delicate nature.

He seems to have been born with 3 skins too few.

Unlike Einstein, who had to subjugate his powerful ego in the study of the external world before he could attain his moral stature, Hardy had to strengthen an ego which wasn't much protected.

This at times in later life made him self-assertive (as Einstein never was) when he had to take a moral stand.

On the other hand, it gave him his introspective insight and beautiful candor, so that he could speak of himself with absolute simplicity (as Einstein never could).

%
I believe this contradiction, or tension, in this temperament was linked with a curious tic in his behavior.

He was the classical anti-narcissist.

He could not endure having his photograph taken: so far as I know, there are only 5 snapshots in existence.

He would not have any looking glass in his rooms, not even a shaving mirror.

When he went to a hotel, his 1st action was to cover all the looking-glass with towels.

This would have been odd enough, if his face had been like a gargoyle: superficially it might seem odder, since all his life he was good-looking quite out of the ordinary.

But, of course, narcissism and anti-narcissism have nothing to do with looks as outside observers see them.

%
This behavior seems eccentric, and indeed it was.

Between him and Einstein, though, there was a difference in kind.

Those who spent much time with Einstein - e.g. Infeld - found him grow stranger, less like themselves, the longer they knew him.

I am certain that I should have felt the same.

With Hardy the opposite was true.

His behavior was often different, bizarrely so, from ours: but it came to seem a kind of superstructure set upon a nature which wasn't all that different from our own, except that it was more delicate, less padded, finer-nerved.

%
The other unusual feature of his childhood was more mundane: but it meant the removal of all practical obstacles throughout his entire career.

Hardy, with his limpid honesty, would have been the last man to be finicky on this matter.

He knew what privilege meant, and he knew that he had possessed it.

His family had no money, only a schoolmaster's income, but they were in touch with the best educational advice of late 19th-century England.

That particular kind of information has always been more significant in this country than any amount of wealth.

The scholarships have been there all right, if one knew how to win them.

There was never the slightest chance of the young Wells or the young Einstein.

From the age of 12 he had only to survive, and his talents would be looked after.

%
At 12, in fact, he was given a scholarship at Winchester, then and for long afterwards the best mathematical school in England, simply on the strength of some mathematical work he had done at Granleigh.

(Incidentally, one wonders if any great school could be so elastic nowadays?)

There he was taught mathematics in a class of one: in classics he was as good as the other top collegers.

Later, he admitted that he had been well-educated, but he admitted it reluctantly.

He disliked the school, except for its classes.

Like all Victorian public schools, Winchester was a pretty rough place.

He nearly died 1 winter.

He envied Littlewood in his cared-for home as a day boy at St Paul's or other friends at our free-and-easy grammar schools.

He never went near Winchester after he had left it: but he left it, with the inevitability of one who had got on to the right tramlines, with an open scholarship to Trinity.

%
He had 1 curious grievance against Winchester.

He was a natural ball-games player with a splendid eye.

In this 50s he could usually beat the university 2nd string at real tennis, and in his 60s I saw him bring off startling shots in the cricket nets.

Yet he had not had an hour's coaching at Winchester: his method was defective: if he had been coached, he thought, he would have been a really good batsman, not quite 1st-class, but not too far away.

Like all his judgments on himself, I believe that one is quite true.

It is strange that, at the zenith of Victorian games-worship, such a talent was utterly missed.

I supposed no one thought it worth looking for in the school's top scholar, so frail and sickly, so defensively shy.

%
It would have been natural for a Wykehamist of his period to go to New College.

That wouldn't have made much difference to his professional career (though, since he always liked Oxford better than Cambridge, he might have stayed there all his life, and some of us would have missed a treat).

He decided to go to Trinity instead, for a reason that he describes, humorously but with his usual undecorated truth, in the \textit{Apology}.
\begin{quotation}
    `I was about 15 when (in a rather odd way) my ambitions took a sharper turn.
    
    These is a book by ``Alan St Aubyn'' (actually Mrs Frances Marshall) called \textit{A Fellow of Trinity}, 1 of a series dealing with what is supposed to be Cambridge college life$\ldots$
    
    There are 2 heroes, a primary hero called Flowers, who is almost wholly good, and a secondary hero, a much weaker vessel, called Brown.
    
    Flowers and Brown find many dangers in university life$\ldots$
    
    Flowers survives all these troubles, is Second Wrangler and succeeds automatically to a Fellowship (as I suppose he would have done then).
    
    Brown succumbs, ruins his parents, takes to drink, is saved from delirium tremens during a thunderstorm only by the prayers of the Junior Dean, has much difficulty in obtaining even an Ordinary Degree, and ultimately becomes a missionary.
    
    The friendship is not shattered by these unhappy events, and Flowers's thoughts stray to Brown, with affectionate pity, as he drinks port and eats walnuts for the 1st time in Senior Combination Room.
    
    `Now Flowers was a decent enough fellow (so far as ``Alan St Aubyn'' could draw one), but even my unsophisticated mind refused to accept him as clever.
    
    If he could do these things, why not I?
    
    In particular, the final scene in Combination Room fascinated me completely, and from that time, until I obtained one, mathematics meant to me primarily a Fellowship of Trinity.'
\end{quotation}
Which he duly obtained, after getting the highest place in the Mathematical Tripos Part II, at the age of 22.

On the way, there were 2 minor vicissitudes.

The 1st was theological, in the high Victorian manner.

Hardy had decided - I think before he left Winchester - that he did not believe in God.

With him, this was a black-and-white decision, as sharp and clear as all other concepts in his mind.

Chapel at Trinity was compulsory.

Hardy told the Dean, no doubt with his own kind of shy certainty, that he could not conscientiously attend.

The Dean, who must have been a jack-in-office, insisted that Hardy should write to his parents and tell them so.

They were orthodox religious people, and the Dean knew, and Hardy knew much more, that the news would give them pain - pain such as we, 70 years later, cannot easily imagine.

%
Hardy struggled with his conscience.

He wasn't worldly enough to slip the issue.

He wasn't even worldly enough - he told me 1 afternoon at Fenner's, for the wound still rankled - to take the advice of more sophisticated friends, e.g. George Trevelyan and Desmond MacCarthy, who would have known how to handle the matter.

In the end he wrote the letter.

Partly because of that incident, his religious disbelief remained open and active ever after.

He refused to go into any college chapel even for formal business, like electing a master.

He had clerical friends, but God was his personal enemy.

In all this there was an echo of the 19th century: but one would be wrong, as always with Hardy, not to take him at his word.

%
Still, he turned it into high jinks.

I remember, 1 day in the 30s, seeing him enjoy a minor triumph.

It happened in a Gentleman $v$.

Players match at Lord's.

It was early in the morning's play, and the sun was shining over the pavilion.

1 of the batsmen, facing the Nursery end, complained that he was unsighted by a reflection from somewhere unknown.

The umpires, puzzled, padded round by the sight-screen.

Motor-cars?

No.

Windows?

None on that side of the ground.

At last, with justifiable triumph, an umpire traced the reflection down - it came from a large pectoral cross reposing on the middle of an enormous clergyman.

Politely the umpire asked him to take it off.

Close by, Hardy was doubled up in Mephistophelian delight.

That lunch time, he had no leisure for eating: he was writing postcards (postcards and telegrams were his favorite means of communication) to each of his clerical friends.

%
But in his war against God and God's surrogates, victory was not all on 1 side.

On a quiet and lovely May evening at Fenner's, round about the same period, the chimes of 6 o'clock fell across the ground.

`It's rather unfortunate,' said Hardy simply, `that some of the happiest hours of my life should have been spent within sound of a Roman Catholic church.'

%
The 2nd minor upset of his undergraduate years was professional.

Almost since the time of Newton, and all through the 19th century, Cambridge had been dominated by the examination for the old Mathematical Tripos.

The English have always had more faith in competitive examinations than any other people (except perhaps in Imperical Chinese): they have conducted these examinations with traditional justice: but they have often shown remarkable woodenness in deciding what the examinations should be like.

I.e., incidentally, true to this day.

It was certainly true of the Mathematical Tripos in its glory.

It was an examination in which the questions were usually of considerable mechanical difficulty - but unfortunately did not give any opportunity for the candidate to show mathematical imagination or insight or any quality that a creative mathematicians needs.

The top candidates (the Wranglers - a term which still survives, meaning a 1st Class) were arranged, on the basis of marks, in strict numerical order.

Colleges had celebrations when 1 of their number became Senior Wrangler: the 1st 2 or 3 Wranglers were immediately elected Fellows.

%
It was all very English.

It had only 1 disadvantage, as Hardy pointed out with his polemic clarity, as soon as he had become an eminent mathematician and was engaged, together with his tough ally Littlewood, in getting the system abolished: it had effectively ruined serious mathematics in English for 100 years.

%
In his 1st term at Trinity, Hardy found himself caught in this system.

He ws to be trained as a racehorse, over a course of mathematical exercises which at 19 he knew to be meaningless.

He was sent to a famous coach, to whom most potential Senior Wranglers went.

This coach knew all the obstacles, all the tricks of the examiners, and was sublimely uninterested in the subject itself.

At this point the young Einstein would have rebelled: he would either have left Cambridge or done no formal work for the next 3 years.

But Hardy was born into the more intensely professional English climate (which was its merits as well as its demerits).

After considering changing his subject to history, he had the sense to find a real mathematician to teach him.

Hardy paid him a tribute in the \textit{Apology}:
\begin{quotation}
    `My eyes were 1st opened by Prof. Love, who taught me for a few terms and gave me my 1st serious conception of analysis.
    
    But the great debt which I owe to him - he was, after all, primarily an applied mathematician - was his advice to read Jordan's famous \textit{Cours d' Analyse}: and I shall never forget the astonishment with which I read that remarkable work, the 1st inspiration for so many mathematicians of my generation, and learned for the 1st time as I read it what mathematics really meant.
    
    From that time onwards I was in my way a real mathematician, with sound mathematical ambitions and a genuine passion for mathematics.'
\end{quotation}
He was 4th Wrangler in 1898.

This faintly irritated him, he used to confess.

He was enough of a natural competitor to feel that, though the race was ridiculous, he ought to have won it.

In 1900 he took Part II of the Tripos, a more respect-worthy examination, and got his right place and his Fellowship.

%
From that time on, his life was in essence settled.

He knew his purpose, which was to bring rigor into English mathematical analysis.

He did not deviate from the researches which he called `the 1 great permanent happiness of my life'.

There were no anxieties about what he should do.

Neither he nor anyone else doubted his great talent.

He was elected to the Royal Society at 33.

%
In many senses, then, he was unusually lucky.

He did not have to think about his career.

From the time he was 23 he had all the leisure that a man could want, and as much money as he needed.

A bachelor don in Trinity in the 1900's was comfortably off.

Hardy was sensible about money, spent it when he felt impelled (sometimes for singular purposes, e.g. 50-mile taxi-rides), and otherwise was not at all unworldly about investments.

He played his games and indulged hi eccentricities.

He was living in some of the best intellectual company in the world - G. E. Moore, Whitehead, Bertrand Russel, Trevelyan, the high Trinity society which was shortly to find its artistic complement in Bloomsbury.

(Hardy himself had links with Bloomsbury, both of personal friendship and of sympathy.)

In a brilliant circle, he was 1 of the most brilliant young men - and, in a quiet way, 1 of the most irrepressible.

%
I will anticipate now what I shall say later.

His life remained the life of a brilliant young man until he was old: so did his spirit: his games, his interests, kept the lightness of a young don's.

And, like many men who keep a young man's interests into their 60s, his last years were the darker for it.

%
Much of his life, though, he was happier than most of us.

He had a great many friends, of surprising different kinds.

These friends had to pass some of his private tests: they needed to possess a quality which he called `spin' (this is a cricket term, and untranslatable: it implies a certain obliquity for irony of approach: of recent public figures, Macmillan and Kennedy would get high marks for spin, Churchill and Eisenhower not).

But he was tolerant, loyal, extremely high-spirited, and in an undemonstrative way found of his friends.

I once was compelled to go and see him in the morning, which was always his set time for mathematical work.

He was sitting at his desk, writing in his beautiful calligraphy.

I murmured some commonplace politeness about hoping that I wasn't disturbing him.

He suddenly dissolved into his mischievous grin.

`As you ought to be able to notice, the answer to that is that you are.

Still, I'm usually glad to see you.'

In the 16 years we knew each other, he didn't say anything more demonstrative than that: except on his deathbed, when he told me that he looked forward to my visits.

%
I think my experience was shared by most of his close friends.

But he had, scattered through his life, 2 or 3 other relationships, different of kind.

These were intense affections, absorbing, non-physical but exalted.

The one I knew about was for a young man whose nature was as spiritually delicate as his own.

I believe, though I only picked this up from chance remarks, that the same was true of the others.

To many people of my generation, such relationships would seem either unsatisfactory or impossible.

They were neither the one nor the other; and, unless one takes them for granted, one doesn't begin to understand the temperament of men like Hardy (they are rare, but not as rare as white rhinoceroses), nor the Cambridge society of his time.

He didn't get the satisfactions that most of us can't help finding: but he knew himself unusually well, and that didn't make him unhappy.

His inner life was his own, and very rich.

The sadness came at the end.

Apart from his devoted sister, he was left with no one close to him.

%


''
\part{}

\chapter{}

\section{}


%------------------------------------------------------------------------------%

\printbibliography[heading=bibintoc]
\end{document}